\documentclass[12pt]{article}
\usepackage[utf8]{inputenc}
\usepackage[T1]{fontenc}
\usepackage{hyperref}
\title{Zaza (Zazaki) Language Resources for K-12 Curriculum Development}
\author{Compiled by Claude Sonnet 5.5 (claude-sonnet-5-5) for Andr\'as Kornai}
\date{}
\begin{document}
\maketitle

\section{Introduction}
The prompt gave only the code zza, which is the ISO 639-3 code for Zaza (Zazaki), an Indo-European language of eastern Turkey. This report treats zza as Zaza. The search did not confirm the code itself. Several teaching and reference resources exist, but the search found no school curriculum or programme for Zazaki.

\section{Primary Educational Resources}
\begin{itemize}
\item Mesut Keskin, \emph{Zon\^e Ma Zanena -- Zazaki for Beginners} (2nd edition, 2016), a beginner textbook. It is known here only through its companion word list on Webonary, so this report does not give a publisher link for the textbook.
\item Digital Dialects: free illustrated picture dictionaries for beginner learners. \url{https://digitaldialects.com/Zazaki_picture_dictionaries.htm}
\end{itemize}

\section{Linguistic and Reference Materials}
\begin{itemize}
\item Webonary Zazaki dictionary, built from a Zazaki-German and German-Zazaki word list. \url{https://www.webonary.org/zazaki/foreword/}
\item Entry listing from the same word list. \url{https://www.webonary.work/zazaki/?p=184575}
\item Nicholas Awde, \emph{Zaza-English / English-Zaza Dictionary \& Phrasebook} (Hippocrene Books). Listings disagree on its size (3,000 versus over 8,500 entries). \url{https://www.hippocrenebooks.com/store/p513/ZazaDictionaryandPhrasebook.html}
\item Munzur University course content, Zaza Language and Literature. \url{https://www.munzur.edu.tr/birimler/idari/erasmus/pages/File/DersIcerikleri/Faculty%20of%20Letters/Zaza%20Language%20and%20Literature.pdf}
\end{itemize}

\section{Religious and Community Resources}
The search found no religious or community resources for Zazaki.

\section{Digital and Audio Resources}
The Webonary dictionary and the Digital Dialects picture dictionaries above are the digital resources found. The search found no separate audio resources.

\section{Curriculum Development Context}
Zazaki has university-level course content (Munzur University, Turkey), a beginner textbook and a few dictionaries. It has no school curriculum or programme in these results. A K-12 programme would need to build on the Keskin textbook and the Webonary lexicon, and would need new graded materials for each school level.

\end{document}